\chapter*{Glossary}
\addcontentsline{toc}{chapter}{Glossary}
These entries prevent distinct model and runtime objects from being conflated.
They describe the manuscript's interfaces, not demonstrated biological capabilities.
\begin{description}
\item[DOGMA] Proposed non-Transformer DNA-native model/runtime research direction.
\item[Hermon DNA] Transformer-based DNA/genomic model/runtime direction.
\item[Evolutor] Broader orchestration, research, compiler/runtime, memory,
adaptation and governance layer; not the external model named Evo 2.
\item[Parameter] Learned value shared across uses of a model operation.
\item[Activation state] Values produced for a particular input or continuing
record with fixed parameters; not automatically learned persistent knowledge.
\item[Optimizer state] Learning-procedure memory such as moments or momentum,
distinct from a model's record-level carry.
\item[Carried state] All fields needed to resume the same computation under
the same configuration, including boundary metadata.
\item[Checkpoint] Version-bound stored snapshot of the declared model or
learning/runtime state; a parameter file alone may not restore the full run.
\item[Padding identity] Storage padding leaves the logical state unchanged.
An uncertain biological base is not automatically padding.
\item[Record reset] Explicit replacement with initial state at an independent
record boundary; neither zero retention nor a chunk boundary implies reset.
\item[Update clock] Events, elapsed time or genomic-distance increments that
advance an explicitly defined state update.
\item[Retention] Attenuation of old state in a specified recurrence; not proof
of task-useful long-range memory.
\item[Memory read] Access to retained information under a declared addressing
and information-availability policy.
\item[Equivariance] Input transformation induces a specified output or feature
transformation; this differs from leaving the result unchanged.
\item[Invariance] A result remains unchanged under a declared transformation.
Strand-sensitive labels may require a nontrivial label action instead.
\item[Offline/full-record access] Available information can include suffix
content. The Chapter 21 dual/RC representation belongs to this contract.
\item[Causal feature] Depends only on the permitted prefix and declared
external knowledge. Target alignment and token spans must also be legal.
\item[Gradient parity] Agreement of derivatives as well as output values;
detaching chunk state can preserve one and break the other.
\item[Functional ablation] Intervention on one mechanism with a declared fixed
parameter set; retraining a controlled alternative answers a different question.
\item[Non-identifiability] Distinct parameter explanations implement the same
observed function under the declared measurement interface.
\item[Token span] Half-open interval in nucleotide coordinates carried by a token.
A token ordinal is not its physical base coordinate or width.
\item[Joint strand action] Sequence reverse complement together with the
specified orientation/feature/label transformation.
\item[Rotary position] Pairwise feature rotation indexed by declared position
units; relative-score algebra does not prove unseen-context generalization.
\item[Legal attention edge] A query-key access allowed by all mask, validity,
information and objective constraints.
\item[Global relay] A designated within-record attention position aggregating
information; another layer reads its updated representation.
\item[Retrieval] Access to a specified external corpus under a legal query and
knowledge policy, not merely a global attention token.
\item[Capability evidence] Task- and protocol-bound observations of performance;
implemented memory or agent persistence alone is not evidence of general intelligence.
\item[Estimand] Precisely defined quantity an experiment intends to estimate,
including predictors or training procedure, target population and sampling unit.
\item[Dependency closure] Connected components of all declared relationships
that prohibit cross-partition placement; curation adequacy remains separate.
\item[Paired comparison] Evaluate both predictors on the same record and retain
their joint outcomes when estimating differences and uncertainty.
\item[Cluster resampling] Draw independent clusters with replacement while
retaining their dependent records and paired predictions within each draw.
\item[Conditional test uncertainty] Evaluation-sampling variation for fixed
predictors, not automatically variation across training procedures or seeds.
\item[Resource inventory] Static auditable counts and units, distinct from
instrumented workload measurements such as peak memory, latency or energy.
\item[Affine prefix summary] Chronological pair $(A,B)$ mapping incoming state
to $A\odot h+B$; each prefix is required, not merely a terminal reduction.
\item[Scan eligibility] Coefficients must be available without computing the
recurrent states to be scanned; causality alone does not establish eligibility.
\item[Reset-before-update] Replace incoming state by a declared reset state,
then apply the current valid transition; distinct from identity padding.
\item[Boundary detach] Preserve a carried value while intentionally severing its
earlier derivative path; not a transparent gradient-preserving optimization.
\item[Parallel dependency depth] Idealized dependent stage count, distinct from
total work, memory traffic and measured wall-clock latency.
\item[Global weighted reduction] Sum all weighted valid-element loss numerators
and divide by their global positive mass, not an unweighted mean of local means.
\item[Gradient accumulation] Add microbatch derivative contributions at fixed
parameters before a single complete optimizer update.
\item[Activation checkpointing] Recompute omitted forward intermediates during
backward under the same declared input, parameter and stochastic-state contract.
\item[Stochastic replay] Reproduce the required random state during recomputation;
matching first-forward loss alone does not prove gradient equivalence.
\item[Rank mass] Sum of fixed valid-element weights on a rank; not necessarily its record count.
\item[Tensor partition] Disjoint operand/output coordinate ownership with a derived sum or concatenation contract.
\item[Persistent shard] Owned stored array range; does not exclude transient full materialization.
\item[Artifact integrity] Byte/schema agreement relative to a declared identity; not independently authenticated authorship.
\item[Trusted pin] Content identity obtained independently of the artifact being checked.
\item[Inference reload] Reconstruct and evaluate a declared model from artifact bytes without reusing its live training object.
\item[Paired percentile interval] Interval under the stated empirical resampling law for shared test cases; it does not include all training or selection uncertainty automatically.
\item[Crossed bootstrap] Independent resampling of seed and group axes, with the same drawn groups shared across seed rows for a declared equal-axis estimand.
\item[Experimental unit] Independently assigned or sampled unit under the design; duplicated windows do not create new units.
\item[Benchmark envelope] Shared records, target, legal access, metric and opportunity, while preserving architecture-native state and cost.
\item[Linearization point] Declared serialized instant at which an operation takes effect relative to other operations.
\item[ABA reuse] A familiar identifier returns after release; increasing generations distinguish old leases from the new owner.
\item[Terminal reason] Immutable lifecycle outcome, distinct from output draining and physical resource cleanup.
\item[Quiescence] No pending physical work may access the allocation; logical cancellation alone does not establish it.
\item[Backpressure] Refusal to launch additional work when the declared output reservation cannot fit; not silent output dropping.
\end{description}
